% Insert into the manuscript with % Insert into the manuscript with % Insert into the manuscript with % Insert into the manuscript with \input{paper/measurement_architecture}.
% No extra LaTeX packages are required by this fragment.
\subsection{Comparison Architecture and Timing Boundary}
\label{subsec:laptop_timing}
A laptop running ROS~2 Jazzy issues every experimental command and records
all timing events using its local monotonic clock. For the ROX-Diff, the
native path is laptop--ROS~2/DDS--Nav2 on ROX. The standardized path is
laptop--MQTT broker on the Raspberry Pi--VDA~5050 adapter on ROX--the same
Nav2 action server. For the crane, the native laptop accesses the crane's
OPC~UA server directly; the standardized path passes through the Pi broker
and the crane adapter on the Pi before accessing that same OPC~UA server.
The Pi retains its existing crane watchdog in both conditions. The normal
master remains idle; its dashboard may display copies of the laptop's events.
No ROS client container or SSH command launcher is part of these paths.

For each trial, the laptop records command submission $t_0$, acknowledgement
receipt $t_{\mathrm{ack}}$, and completion receipt $t_{\mathrm{completion}}$:
\begin{equation}
T_{\mathrm{ack}}=t_{\mathrm{ack}}-t_0,\qquad
T_{\mathrm{completion}}=t_{\mathrm{completion}}-t_0.
\end{equation}
Connections, payload construction, validation and
start-position checks occur before $t_0$; submission, serialization and
communication occur within the measured interval. ROS and MQTT response
callbacks capture time on entry, before decoding or logging their payload.
The native crane client records the corresponding local API return. No
remote timestamp is subtracted and synchronized device clocks are unnecessary.

For ROX, acknowledgement means the actual Nav2 goal response, and completion
means its terminal action result. In the VDA condition, the adapter forwards
these responses immediately on a dedicated experimental MQTT feedback topic,
with order identifiers and the Nav2 goal UUID. This feedback is instrumentation,
not a standard VDA~5050 acknowledgement, and its serialization and return
transport are included. Periodic VDA state publication is not used as a proxy
for Nav2 acceptance or completion.

The crane executes the same bounded hoist routine in both conditions: directly
on the laptop for native operation and inside the Pi adapter for VDA operation.
Acknowledgement follows successful return of the first OPC~UA motion-write
batch, not local target assignment. Completion follows target attainment,
STOP, and a configured interval of stationary feedback. This is an
application-level completion criterion, not a PLC action-result message.
A benchmark reservation prevents overlapping adapter commands; a heartbeat
lease cancels the experiment if the laptop disappears. The existing watchdog,
manual enable and physical safety controls remain independent. These differing
ROX and crane acknowledgement semantics are reported separately.

Matched pairs use the same verified starting endpoint, target, motion routine,
controller settings and network arrangement, randomizing native/VDA order.
A return-to-start movement between treatments is logged as setup and excluded
from timing summaries. Independent endpoint verification follows the measured
completion event; failures and timeouts remain in the dataset. For either
metric $k\in\{\mathrm{ack},\mathrm{completion}\}$, the paired difference is
\begin{equation}
\Delta T_{k,i}=T_{k,i}^{\mathrm{VDA}}-T_{k,i}^{\mathrm{native}}.
\end{equation}
Report successful/scheduled counts, mean and sample standard deviation,
median and interquartile range, and paired effects with bootstrap confidence
intervals. Upper-tail estimates from small pilot samples are descriptive only.
The comparison measures deployed architecture response times, including
transport, adapter queuing, processing and feedback. For the crane it also
includes the placement of its OPC~UA polling loop on the laptop versus Pi.
It does not isolate adapter CPU cost, one-way latency or pure physical motion
duration, and does not establish complete-handover throughput.

% Executable setup notes (not printed in the paper):
% Apply the patch on laptop, Pi and ROX. Rebuild the ROX adapter on ROX:
%   cd ros2_ws
%   colcon build --packages-select rox_vda5050_adapter --symlink-install
% Restart its usual launch command after sourcing install/setup.bash.
% Restart the Pi master to display schema-3 laptop events. Keep missions idle.
%
% Laptop, repository root (use the same ROS_DOMAIN_ID/RMW as ROX):
%   source /opt/ros/jazzy/setup.bash
%   python3 -m venv --system-site-packages .venv-benchmark
%   source .venv-benchmark/bin/activate
%   python3 -m pip install -r benchmark/requirements.txt
%   python3 -m unittest discover -s tests -v
%   python3 benchmark/logger_overhead.py --output logger-overhead.json
% Set benchmark/config/rox_benchmark.yaml runner_hostname from `hostname`,
% mqtt.host to Pi, and verified map/pose/headings; then configured: true.
% Use a NEW unique run-directory name every time (old clock data incompatible).
%   python3 benchmark/latency_benchmark.py prepare --config benchmark/config/rox_benchmark.yaml --run-dir runs/rox-pilot-laptop --pairs 5 --seed 5050
%   python3 benchmark/latency_benchmark.py check --run-dir runs/rox-pilot-laptop
%   python3 benchmark/latency_benchmark.py run --run-dir runs/rox-pilot-laptop --end-row 1
% Inspect the physical endpoint and log before continuing under supervision:
%   python3 benchmark/latency_benchmark.py run --run-dir runs/rox-pilot-laptop
% The next unattempted row is selected automatically; failed attempts stop a campaign.
% Five pairs generate 10 measured movements + 5 excluded reset movements.
%
% Crane: verify benchmark/config/crane_benchmark.yaml values at Ilmatar first.
% Copy that exact configured file to Pi; in the crane adapter service environment:
%   CRANE_BENCHMARK_CONFIG=/absolute/path/to/benchmark/config/crane_benchmark.yaml
% Restart the existing crane adapter/watchdog service with that environment.
% Leave its existing URL/access-code/watchdog settings intact. The laptop needs
% direct network access to the PLC URL; it never sends watchdog/access-code writes.
% Unset CRANE_BENCHMARK_CONFIG and restart after tests to restore ordinary orders.
% Prepare/check/run using the SAME CLI with --config benchmark/config/crane_benchmark.yaml
% and a distinct --run-dir runs/crane-pilot-laptop. No ROS installation needed on Pi.
%
% Analysis, on laptop, for either campaign (replace run directory as appropriate):
%   python3 analysis/derive_latency.py runs/rox-pilot-laptop/events.jsonl --schedule runs/rox-pilot-laptop/schedule.csv --output runs/rox-pilot-laptop/trials.csv
%   python3 analysis/statistics.py runs/rox-pilot-laptop/trials.csv --output-dir runs/rox-pilot-laptop/results
% Outputs: trials.csv, latency_summary.csv, paired_differences.csv,
% trial_accounting.csv, latency_summary_table.tex, latency_comparison.pdf/png.
% Freeze environment evidence before the final campaign:
%   python3 benchmark/freeze_environment.py --config benchmark/config/rox_benchmark.yaml --output runs/rox-pilot-laptop/environment.json --network-note "Laptop to Pi broker and ROX/PLC; record wired/Wi-Fi links"
% Run freeze_environment.py on Pi too to capture its broker/crane environment;
% retain both files with the campaign. Record PLC version and Nav2 configuration.
% After pilot review, prepare a NEW run with --pairs 30 and a prespecified seed.
% 30 pairs generate 60 measured movements + 30 excluded resets per device.
.
% No extra LaTeX packages are required by this fragment.
\subsection{Comparison Architecture and Timing Boundary}
\label{subsec:laptop_timing}
A laptop running ROS~2 Jazzy issues every experimental command and records
all timing events using its local monotonic clock. For the ROX-Diff, the
native path is laptop--ROS~2/DDS--Nav2 on ROX. The standardized path is
laptop--MQTT broker on the Raspberry Pi--VDA~5050 adapter on ROX--the same
Nav2 action server. For the crane, the native laptop accesses the crane's
OPC~UA server directly; the standardized path passes through the Pi broker
and the crane adapter on the Pi before accessing that same OPC~UA server.
The Pi retains its existing crane watchdog in both conditions. The normal
master remains idle; its dashboard may display copies of the laptop's events.
No ROS client container or SSH command launcher is part of these paths.

For each trial, the laptop records command submission $t_0$, acknowledgement
receipt $t_{\mathrm{ack}}$, and completion receipt $t_{\mathrm{completion}}$:
\begin{equation}
T_{\mathrm{ack}}=t_{\mathrm{ack}}-t_0,\qquad
T_{\mathrm{completion}}=t_{\mathrm{completion}}-t_0.
\end{equation}
Connections, payload construction, validation and
start-position checks occur before $t_0$; submission, serialization and
communication occur within the measured interval. ROS and MQTT response
callbacks capture time on entry, before decoding or logging their payload.
The native crane client records the corresponding local API return. No
remote timestamp is subtracted and synchronized device clocks are unnecessary.

For ROX, acknowledgement means the actual Nav2 goal response, and completion
means its terminal action result. In the VDA condition, the adapter forwards
these responses immediately on a dedicated experimental MQTT feedback topic,
with order identifiers and the Nav2 goal UUID. This feedback is instrumentation,
not a standard VDA~5050 acknowledgement, and its serialization and return
transport are included. Periodic VDA state publication is not used as a proxy
for Nav2 acceptance or completion.

The crane executes the same bounded hoist routine in both conditions: directly
on the laptop for native operation and inside the Pi adapter for VDA operation.
Acknowledgement follows successful return of the first OPC~UA motion-write
batch, not local target assignment. Completion follows target attainment,
STOP, and a configured interval of stationary feedback. This is an
application-level completion criterion, not a PLC action-result message.
A benchmark reservation prevents overlapping adapter commands; a heartbeat
lease cancels the experiment if the laptop disappears. The existing watchdog,
manual enable and physical safety controls remain independent. These differing
ROX and crane acknowledgement semantics are reported separately.

Matched pairs use the same verified starting endpoint, target, motion routine,
controller settings and network arrangement, randomizing native/VDA order.
A return-to-start movement between treatments is logged as setup and excluded
from timing summaries. Independent endpoint verification follows the measured
completion event; failures and timeouts remain in the dataset. For either
metric $k\in\{\mathrm{ack},\mathrm{completion}\}$, the paired difference is
\begin{equation}
\Delta T_{k,i}=T_{k,i}^{\mathrm{VDA}}-T_{k,i}^{\mathrm{native}}.
\end{equation}
Report successful/scheduled counts, mean and sample standard deviation,
median and interquartile range, and paired effects with bootstrap confidence
intervals. Upper-tail estimates from small pilot samples are descriptive only.
The comparison measures deployed architecture response times, including
transport, adapter queuing, processing and feedback. For the crane it also
includes the placement of its OPC~UA polling loop on the laptop versus Pi.
It does not isolate adapter CPU cost, one-way latency or pure physical motion
duration, and does not establish complete-handover throughput.

% Executable setup notes (not printed in the paper):
% Apply the patch on laptop, Pi and ROX. Rebuild the ROX adapter on ROX:
%   cd ros2_ws
%   colcon build --packages-select rox_vda5050_adapter --symlink-install
% Restart its usual launch command after sourcing install/setup.bash.
% Restart the Pi master to display schema-3 laptop events. Keep missions idle.
%
% Laptop, repository root (use the same ROS_DOMAIN_ID/RMW as ROX):
%   source /opt/ros/jazzy/setup.bash
%   python3 -m venv --system-site-packages .venv-benchmark
%   source .venv-benchmark/bin/activate
%   python3 -m pip install -r benchmark/requirements.txt
%   python3 -m unittest discover -s tests -v
%   python3 benchmark/logger_overhead.py --output logger-overhead.json
% Set benchmark/config/rox_benchmark.yaml runner_hostname from `hostname`,
% mqtt.host to Pi, and verified map/pose/headings; then configured: true.
% Use a NEW unique run-directory name every time (old clock data incompatible).
%   python3 benchmark/latency_benchmark.py prepare --config benchmark/config/rox_benchmark.yaml --run-dir runs/rox-pilot-laptop --pairs 5 --seed 5050
%   python3 benchmark/latency_benchmark.py check --run-dir runs/rox-pilot-laptop
%   python3 benchmark/latency_benchmark.py run --run-dir runs/rox-pilot-laptop --end-row 1
% Inspect the physical endpoint and log before continuing under supervision:
%   python3 benchmark/latency_benchmark.py run --run-dir runs/rox-pilot-laptop
% The next unattempted row is selected automatically; failed attempts stop a campaign.
% Five pairs generate 10 measured movements + 5 excluded reset movements.
%
% Crane: verify benchmark/config/crane_benchmark.yaml values at Ilmatar first.
% Copy that exact configured file to Pi; in the crane adapter service environment:
%   CRANE_BENCHMARK_CONFIG=/absolute/path/to/benchmark/config/crane_benchmark.yaml
% Restart the existing crane adapter/watchdog service with that environment.
% Leave its existing URL/access-code/watchdog settings intact. The laptop needs
% direct network access to the PLC URL; it never sends watchdog/access-code writes.
% Unset CRANE_BENCHMARK_CONFIG and restart after tests to restore ordinary orders.
% Prepare/check/run using the SAME CLI with --config benchmark/config/crane_benchmark.yaml
% and a distinct --run-dir runs/crane-pilot-laptop. No ROS installation needed on Pi.
%
% Analysis, on laptop, for either campaign (replace run directory as appropriate):
%   python3 analysis/derive_latency.py runs/rox-pilot-laptop/events.jsonl --schedule runs/rox-pilot-laptop/schedule.csv --output runs/rox-pilot-laptop/trials.csv
%   python3 analysis/statistics.py runs/rox-pilot-laptop/trials.csv --output-dir runs/rox-pilot-laptop/results
% Outputs: trials.csv, latency_summary.csv, paired_differences.csv,
% trial_accounting.csv, latency_summary_table.tex, latency_comparison.pdf/png.
% Freeze environment evidence before the final campaign:
%   python3 benchmark/freeze_environment.py --config benchmark/config/rox_benchmark.yaml --output runs/rox-pilot-laptop/environment.json --network-note "Laptop to Pi broker and ROX/PLC; record wired/Wi-Fi links"
% Run freeze_environment.py on Pi too to capture its broker/crane environment;
% retain both files with the campaign. Record PLC version and Nav2 configuration.
% After pilot review, prepare a NEW run with --pairs 30 and a prespecified seed.
% 30 pairs generate 60 measured movements + 30 excluded resets per device.
.
% No extra LaTeX packages are required by this fragment.
\subsection{Comparison Architecture and Timing Boundary}
\label{subsec:laptop_timing}
A laptop running ROS~2 Jazzy issues every experimental command and records
all timing events using its local monotonic clock. For the ROX-Diff, the
native path is laptop--ROS~2/DDS--Nav2 on ROX. The standardized path is
laptop--MQTT broker on the Raspberry Pi--VDA~5050 adapter on ROX--the same
Nav2 action server. For the crane, the native laptop accesses the crane's
OPC~UA server directly; the standardized path passes through the Pi broker
and the crane adapter on the Pi before accessing that same OPC~UA server.
The Pi retains its existing crane watchdog in both conditions. The normal
master remains idle; its dashboard may display copies of the laptop's events.
No ROS client container or SSH command launcher is part of these paths.

For each trial, the laptop records command submission $t_0$, acknowledgement
receipt $t_{\mathrm{ack}}$, and completion receipt $t_{\mathrm{completion}}$:
\begin{equation}
T_{\mathrm{ack}}=t_{\mathrm{ack}}-t_0,\qquad
T_{\mathrm{completion}}=t_{\mathrm{completion}}-t_0.
\end{equation}
Connections, payload construction, validation and
start-position checks occur before $t_0$; submission, serialization and
communication occur within the measured interval. ROS and MQTT response
callbacks capture time on entry, before decoding or logging their payload.
The native crane client records the corresponding local API return. No
remote timestamp is subtracted and synchronized device clocks are unnecessary.

For ROX, acknowledgement means the actual Nav2 goal response, and completion
means its terminal action result. In the VDA condition, the adapter forwards
these responses immediately on a dedicated experimental MQTT feedback topic,
with order identifiers and the Nav2 goal UUID. This feedback is instrumentation,
not a standard VDA~5050 acknowledgement, and its serialization and return
transport are included. Periodic VDA state publication is not used as a proxy
for Nav2 acceptance or completion.

The crane executes the same bounded hoist routine in both conditions: directly
on the laptop for native operation and inside the Pi adapter for VDA operation.
Acknowledgement follows successful return of the first OPC~UA motion-write
batch, not local target assignment. Completion follows target attainment,
STOP, and a configured interval of stationary feedback. This is an
application-level completion criterion, not a PLC action-result message.
A benchmark reservation prevents overlapping adapter commands; a heartbeat
lease cancels the experiment if the laptop disappears. The existing watchdog,
manual enable and physical safety controls remain independent. These differing
ROX and crane acknowledgement semantics are reported separately.

Matched pairs use the same verified starting endpoint, target, motion routine,
controller settings and network arrangement, randomizing native/VDA order.
A return-to-start movement between treatments is logged as setup and excluded
from timing summaries. Independent endpoint verification follows the measured
completion event; failures and timeouts remain in the dataset. For either
metric $k\in\{\mathrm{ack},\mathrm{completion}\}$, the paired difference is
\begin{equation}
\Delta T_{k,i}=T_{k,i}^{\mathrm{VDA}}-T_{k,i}^{\mathrm{native}}.
\end{equation}
Report successful/scheduled counts, mean and sample standard deviation,
median and interquartile range, and paired effects with bootstrap confidence
intervals. Upper-tail estimates from small pilot samples are descriptive only.
The comparison measures deployed architecture response times, including
transport, adapter queuing, processing and feedback. For the crane it also
includes the placement of its OPC~UA polling loop on the laptop versus Pi.
It does not isolate adapter CPU cost, one-way latency or pure physical motion
duration, and does not establish complete-handover throughput.

% Executable setup notes (not printed in the paper):
% Apply the patch on laptop, Pi and ROX. Rebuild the ROX adapter on ROX:
%   cd ros2_ws
%   colcon build --packages-select rox_vda5050_adapter --symlink-install
% Restart its usual launch command after sourcing install/setup.bash.
% Restart the Pi master to display schema-3 laptop events. Keep missions idle.
%
% Laptop, repository root (use the same ROS_DOMAIN_ID/RMW as ROX):
%   source /opt/ros/jazzy/setup.bash
%   python3 -m venv --system-site-packages .venv-benchmark
%   source .venv-benchmark/bin/activate
%   python3 -m pip install -r benchmark/requirements.txt
%   python3 -m unittest discover -s tests -v
%   python3 benchmark/logger_overhead.py --output logger-overhead.json
% Set benchmark/config/rox_benchmark.yaml runner_hostname from `hostname`,
% mqtt.host to Pi, and verified map/pose/headings; then configured: true.
% Use a NEW unique run-directory name every time (old clock data incompatible).
%   python3 benchmark/latency_benchmark.py prepare --config benchmark/config/rox_benchmark.yaml --run-dir runs/rox-pilot-laptop --pairs 5 --seed 5050
%   python3 benchmark/latency_benchmark.py check --run-dir runs/rox-pilot-laptop
%   python3 benchmark/latency_benchmark.py run --run-dir runs/rox-pilot-laptop --end-row 1
% Inspect the physical endpoint and log before continuing under supervision:
%   python3 benchmark/latency_benchmark.py run --run-dir runs/rox-pilot-laptop
% The next unattempted row is selected automatically; failed attempts stop a campaign.
% Five pairs generate 10 measured movements + 5 excluded reset movements.
%
% Crane: verify benchmark/config/crane_benchmark.yaml values at Ilmatar first.
% Copy that exact configured file to Pi; in the crane adapter service environment:
%   CRANE_BENCHMARK_CONFIG=/absolute/path/to/benchmark/config/crane_benchmark.yaml
% Restart the existing crane adapter/watchdog service with that environment.
% Leave its existing URL/access-code/watchdog settings intact. The laptop needs
% direct network access to the PLC URL; it never sends watchdog/access-code writes.
% Unset CRANE_BENCHMARK_CONFIG and restart after tests to restore ordinary orders.
% Prepare/check/run using the SAME CLI with --config benchmark/config/crane_benchmark.yaml
% and a distinct --run-dir runs/crane-pilot-laptop. No ROS installation needed on Pi.
%
% Analysis, on laptop, for either campaign (replace run directory as appropriate):
%   python3 analysis/derive_latency.py runs/rox-pilot-laptop/events.jsonl --schedule runs/rox-pilot-laptop/schedule.csv --output runs/rox-pilot-laptop/trials.csv
%   python3 analysis/statistics.py runs/rox-pilot-laptop/trials.csv --output-dir runs/rox-pilot-laptop/results
% Outputs: trials.csv, latency_summary.csv, paired_differences.csv,
% trial_accounting.csv, latency_summary_table.tex, latency_comparison.pdf/png.
% Freeze environment evidence before the final campaign:
%   python3 benchmark/freeze_environment.py --config benchmark/config/rox_benchmark.yaml --output runs/rox-pilot-laptop/environment.json --network-note "Laptop to Pi broker and ROX/PLC; record wired/Wi-Fi links"
% Run freeze_environment.py on Pi too to capture its broker/crane environment;
% retain both files with the campaign. Record PLC version and Nav2 configuration.
% After pilot review, prepare a NEW run with --pairs 30 and a prespecified seed.
% 30 pairs generate 60 measured movements + 30 excluded resets per device.
.
% No extra LaTeX packages are required by this fragment.
\subsection{Comparison Architecture and Timing Boundary}
\label{subsec:laptop_timing}
A laptop running ROS~2 Jazzy issues every experimental command and records
all timing events using its local monotonic clock. For the ROX-Diff, the
native path is laptop--ROS~2/DDS--Nav2 on ROX. The standardized path is
laptop--MQTT broker on the Raspberry Pi--VDA~5050 adapter on ROX--the same
Nav2 action server. For the crane, the native laptop accesses the crane's
OPC~UA server directly; the standardized path passes through the Pi broker
and the crane adapter on the Pi before accessing that same OPC~UA server.
The Pi retains its existing crane watchdog in both conditions. The normal
master remains idle; its dashboard may display copies of the laptop's events.
No ROS client container or SSH command launcher is part of these paths.

For each trial, the laptop records command submission $t_0$, acknowledgement
receipt $t_{\mathrm{ack}}$, and completion receipt $t_{\mathrm{completion}}$:
\begin{equation}
T_{\mathrm{ack}}=t_{\mathrm{ack}}-t_0,\qquad
T_{\mathrm{completion}}=t_{\mathrm{completion}}-t_0.
\end{equation}
Connections, payload construction, validation and
start-position checks occur before $t_0$; submission, serialization and
communication occur within the measured interval. ROS and MQTT response
callbacks capture time on entry, before decoding or logging their payload.
The native crane client records the corresponding local API return. No
remote timestamp is subtracted and synchronized device clocks are unnecessary.

For ROX, acknowledgement means the actual Nav2 goal response, and completion
means its terminal action result. In the VDA condition, the adapter forwards
these responses immediately on a dedicated experimental MQTT feedback topic,
with order identifiers and the Nav2 goal UUID. This feedback is instrumentation,
not a standard VDA~5050 acknowledgement, and its serialization and return
transport are included. Periodic VDA state publication is not used as a proxy
for Nav2 acceptance or completion.

The crane executes the same bounded hoist routine in both conditions: directly
on the laptop for native operation and inside the Pi adapter for VDA operation.
Acknowledgement follows successful return of the first OPC~UA motion-write
batch, not local target assignment. Completion follows target attainment,
STOP, and a configured interval of stationary feedback. This is an
application-level completion criterion, not a PLC action-result message.
A benchmark reservation prevents overlapping adapter commands; a heartbeat
lease cancels the experiment if the laptop disappears. The existing watchdog,
manual enable and physical safety controls remain independent. These differing
ROX and crane acknowledgement semantics are reported separately.

Matched pairs use the same verified starting endpoint, target, motion routine,
controller settings and network arrangement, randomizing native/VDA order.
A return-to-start movement between treatments is logged as setup and excluded
from timing summaries. Independent endpoint verification follows the measured
completion event; failures and timeouts remain in the dataset. For either
metric $k\in\{\mathrm{ack},\mathrm{completion}\}$, the paired difference is
\begin{equation}
\Delta T_{k,i}=T_{k,i}^{\mathrm{VDA}}-T_{k,i}^{\mathrm{native}}.
\end{equation}
Report successful/scheduled counts, mean and sample standard deviation,
median and interquartile range, and paired effects with bootstrap confidence
intervals. Upper-tail estimates from small pilot samples are descriptive only.
The comparison measures deployed architecture response times, including
transport, adapter queuing, processing and feedback. For the crane it also
includes the placement of its OPC~UA polling loop on the laptop versus Pi.
It does not isolate adapter CPU cost, one-way latency or pure physical motion
duration, and does not establish complete-handover throughput.

% Executable setup notes (not printed in the paper):
% Apply the patch on laptop, Pi and ROX. Rebuild the ROX adapter on ROX:
%   cd ros2_ws
%   colcon build --packages-select rox_vda5050_adapter --symlink-install
% Restart its usual launch command after sourcing install/setup.bash.
% Restart the Pi master to display schema-3 laptop events. Keep missions idle.
%
% Laptop, repository root (use the same ROS_DOMAIN_ID/RMW as ROX):
%   source /opt/ros/jazzy/setup.bash
%   python3 -m venv --system-site-packages .venv-benchmark
%   source .venv-benchmark/bin/activate
%   python3 -m pip install -r benchmark/requirements.txt
%   python3 -m unittest discover -s tests -v
%   python3 benchmark/logger_overhead.py --output logger-overhead.json
% Set benchmark/config/rox_benchmark.yaml runner_hostname from `hostname`,
% mqtt.host to Pi, and verified map/pose/headings; then configured: true.
% Use a NEW unique run-directory name every time (old clock data incompatible).
%   python3 benchmark/latency_benchmark.py prepare --config benchmark/config/rox_benchmark.yaml --run-dir runs/rox-pilot-laptop --pairs 5 --seed 5050
%   python3 benchmark/latency_benchmark.py check --run-dir runs/rox-pilot-laptop
%   python3 benchmark/latency_benchmark.py run --run-dir runs/rox-pilot-laptop --end-row 1
% Inspect the physical endpoint and log before continuing under supervision:
%   python3 benchmark/latency_benchmark.py run --run-dir runs/rox-pilot-laptop
% The next unattempted row is selected automatically; failed attempts stop a campaign.
% Five pairs generate 10 measured movements + 5 excluded reset movements.
%
% Crane: verify benchmark/config/crane_benchmark.yaml values at Ilmatar first.
% Copy that exact configured file to Pi; in the crane adapter service environment:
%   CRANE_BENCHMARK_CONFIG=/absolute/path/to/benchmark/config/crane_benchmark.yaml
% Restart the existing crane adapter/watchdog service with that environment.
% Leave its existing URL/access-code/watchdog settings intact. The laptop needs
% direct network access to the PLC URL; it never sends watchdog/access-code writes.
% Unset CRANE_BENCHMARK_CONFIG and restart after tests to restore ordinary orders.
% Prepare/check/run using the SAME CLI with --config benchmark/config/crane_benchmark.yaml
% and a distinct --run-dir runs/crane-pilot-laptop. No ROS installation needed on Pi.
%
% Analysis, on laptop, for either campaign (replace run directory as appropriate):
%   python3 analysis/derive_latency.py runs/rox-pilot-laptop/events.jsonl --schedule runs/rox-pilot-laptop/schedule.csv --output runs/rox-pilot-laptop/trials.csv
%   python3 analysis/statistics.py runs/rox-pilot-laptop/trials.csv --output-dir runs/rox-pilot-laptop/results
% Outputs: trials.csv, latency_summary.csv, paired_differences.csv,
% trial_accounting.csv, latency_summary_table.tex, latency_comparison.pdf/png.
% Freeze environment evidence before the final campaign:
%   python3 benchmark/freeze_environment.py --config benchmark/config/rox_benchmark.yaml --output runs/rox-pilot-laptop/environment.json --network-note "Laptop to Pi broker and ROX/PLC; record wired/Wi-Fi links"
% Run freeze_environment.py on Pi too to capture its broker/crane environment;
% retain both files with the campaign. Record PLC version and Nav2 configuration.
% After pilot review, prepare a NEW run with --pairs 30 and a prespecified seed.
% 30 pairs generate 60 measured movements + 30 excluded resets per device.
